\section{Прикладные особенности}
\label{sec:AppliedAspects}
	\tikzsetfigurename{GraphTheory_Applied_}

В предыдущих разделах мы определили полугруппы, моноиды, группы, полукольца и кольца как <<честные>> алгебраические структуры: сформулировали аксиомы, которым они должны удовлетворять, и привели примеры.
На практике же при реализации алгоритмов анализа графов эти аксиомы часто соблюдаются лишь частично, а на первый план выходят вопросы представления данных, эффективности операций и удобства программного интерфейса.
Данный раздел посвящён именно таким прикладным аспектам~\sidecite{7761646}.

Начнём с фундаментального вопроса~--- представления данных.
При определении матрицы смежности (раздел~\ref{chpt:GraphTheoryIntro}) мы ввели множество $\Opt{L} = \{\Bbbzero\} \cup L$, где $\Bbbzero$ означает <<ребро отсутствует>>.
В языках программирования такой конструкции естественно соответствует тип \texttt{Option} (OCaml, F\#) или \texttt{Maybe} (Haskell): значение \texttt{Some x} означает <<есть ребро с меткой \texttt{x}>>, а \texttt{None}~--- <<ребра нет>>.
Таким образом, матрица смежности графа с метками типа \texttt{Lbl} имеет тип $\texttt{Matrix}\langle\texttt{Option}\langle\texttt{Lbl}\rangle\rangle$.

В реальных графах число рёбер, как правило, существенно меньше $|V|^2$, поэтому матрицы смежности разрежены: подавляющее большинство ячеек содержат \texttt{None}.
Хранить их все нет смысла~--- достаточно хранить лишь те ячейки, которые соответствуют реально существующим рёбрам, то есть имеют значение \texttt{Some x}.
Для этого существует множество форматов разреженных матриц: CSR (Compressed Sparse Row), координатный, Quad Tree и другие.
Выбор формата определяется конкретной задачей и используемой библиотекой, но все они основаны на одной и той же идее: хранить только <<непустые>> ячейки.

Определившись со способом хранения, перейдём к операциям над разреженными матрицами.
Поэлементная операция над двумя матрицами~--- это операция, которая применяется независимо к каждой паре ячеек с одинаковыми индексами.
Наивная типизация такой операции имеет вид
\[
    \texttt{Option}\langle T_1\rangle \to \texttt{Option}\langle T_2\rangle \to \texttt{Option}\langle T_3\rangle.
\]
Однако эта сигнатура недостаточно ограничительна: она позволяет определить операцию, возвращающую \texttt{Some x} даже тогда, когда оба аргумента равны \texttt{None}.
Иными словами, можно <<создать значение из двух пустот>>.
Эта проблема известна в сообществе GraphBLAS как \emph{проблема явных и неявных нулей}~\sidecite{7761646}: разреженная матрица хранит только явные (ненулевые) значения, но комбинация двух не-хранимых нулей не должна порождать новое хранимое значение.

Решение~--- ввести тип $\texttt{AtLeastOne}\langle T_1, T_2\rangle$, который гарантирует, что хотя бы один из аргументов непуст:
\begin{align*}
    \texttt{AtLeastOne}\langle T_1, T_2\rangle =\ &\texttt{Both}\ \texttt{of}\ T_1 * T_2 \\
    \mid\ &\texttt{Left}\ \texttt{of}\ T_1 \\
    \mid\ &\texttt{Right}\ \texttt{of}\ T_2.
\end{align*}
Теперь поэлементная операция получает сигнатуру
\[
    \texttt{op}: \texttt{AtLeastOne}\langle T_1, T_2\rangle \to \texttt{Option}\langle T_3\rangle,
\]
которая на уровне типов гарантирует, что операция никогда не будет вызвана с двумя \texttt{None}.
При этом тип результата остаётся $\texttt{Option}\langle T_3\rangle$: операция всё ещё может вернуть \texttt{None}, и это просто означает, что в данной позиции результирующая ячейка не должна храниться.

Использование $\texttt{AtLeastOne}$ позволяет унифицировать поэлементные операции, заменив несколько специализированных функций одной обобщённой высокопорядковой функцией $\texttt{map2}$~\sidecite{7761646}:
\begin{align*}
    \texttt{map2}\ & : \\
    & \quad \texttt{op}: \texttt{AtLeastOne}\langle T_1, T_2\rangle \to \texttt{Option}\langle T_3\rangle \\
    & \to \texttt{m}_1: \texttt{Matrix}\langle\texttt{Option}\langle T_1\rangle\rangle \\
    & \to \texttt{m}_2: \texttt{Matrix}\langle\texttt{Option}\langle T_2\rangle\rangle \\
    & \to \texttt{result}: \texttt{Matrix}\langle\texttt{Option}\langle T_3\rangle\rangle.
\end{align*}
Одна и та же функция $\texttt{map2}$, параметризованная разными операциями $\texttt{op}$, реализует и поэлементное сложение, и поэлементное умножение, и маскирование.
Рассмотрим конкретные примеры.

Операция поэлементного сложения для целых чисел определяется следующим образом:
\begin{align*}
    \texttt{op\_int\_add}\ (\texttt{Both}\ (x, y)) &= \texttt{Some}\ (x + y), \\
    \texttt{op\_int\_add}\ (\texttt{Left}\ x \mid \texttt{Right}\ y) &= \texttt{Some}\ x.
\end{align*}
(Если в результате сложения получился <<честный>> ноль, можно дополнительно вернуть \texttt{None}, чтобы не сохранять его.)

Операция поэлементного умножения, напротив, возвращает непустое значение только тогда, когда оба аргумента непусты:
\begin{align*}
    \texttt{op\_int\_mult}\ (\texttt{Both}\ (x, y)) &= \texttt{Some}\ (x * y), \\
    \texttt{op\_int\_mult}\ (\texttt{Left}\ x \mid \texttt{Right}\ y) &= \texttt{None}.
\end{align*}

Наконец, маска (оставить значения первой матрицы только там, где во второй матрице есть ненулевые значения) также является частным случаем $\texttt{map2}$:
\begin{align*}
    \texttt{op\_mask}\ (\texttt{Both}\ (x, y)) &= \texttt{Some}\ x, \\
    \texttt{op\_mask}\ (\texttt{Left}\ x \mid \texttt{Right}\ y) &= \texttt{None}.
\end{align*}

Идея параметризации операции может быть распространена и на матричное умножение.
Напомним, что в разделе~\ref{chpt:LinearAlgebraIntro} мы определили матричное умножение над полукольцом $(S, \oplus, \otimes)$ как
\[
    P[i,j] = \bigoplus_{l} M[i,l] \otimes N[l,j],
\]
где $\oplus$~--- операция <<сложения>> (агрегации) частичных результатов, а $\otimes$~--- операция <<умножения>> ячеек.

В терминах типов этому соответствует функция $\texttt{mxm}$:
\begin{align*}
    \texttt{mxm}\ & : \\
    & \quad \texttt{m}_1: \texttt{Matrix}\langle\texttt{Option}\langle T_1\rangle\rangle \\
    & \to \texttt{m}_2: \texttt{Matrix}\langle\texttt{Option}\langle T_2\rangle\rangle \\
    & \to \texttt{op\_mult}: \texttt{Option}\langle T_1\rangle \to \texttt{Option}\langle T_2\rangle \to \texttt{Option}\langle T_3\rangle \\
    & \to \texttt{op\_add}: \texttt{Option}\langle T_3\rangle \to \texttt{Option}\langle T_3\rangle \to \texttt{Option}\langle T_3\rangle \\
    & \to \texttt{zero}: \texttt{Option}\langle T_3\rangle \\
    & \to \texttt{result}: \texttt{Matrix}\langle\texttt{Option}\langle T_3\rangle\rangle.
\end{align*}

Здесь $\texttt{op\_mult}$ задаёт, как перемножить ячейку первой матрицы с ячейкой второй, $\texttt{op\_add}$~--- как сложить несколько частичных произведений для одной позиции результата, а $\texttt{zero}$~--- нейтральный элемент $\texttt{op\_add}$.
Преимущество такого подхода перед классическим определением через полукольцо в том, что $\texttt{op\_mult}$ и $\texttt{op\_add}$ не обязаны удовлетворять аксиомам ассоциативности, коммутативности или дистрибутивности.
На практике при решении задач анализа графов используемые операции часто не образуют полукольцо в математическом смысле, но это не мешает им быть полезными: достаточно, чтобы они были согласованы с конкретной решаемой задачей.
Так, полукольцом может называться структура с доменами разных типов для аргументов умножения, что формально не является полукольцом в алгебраическом смысле, однако такая <<вольность>> оправдана практическими потребностями.

Именно такой подход реализован в стандарте GraphBLAS~\sidecite{7761646}~--- открытом стандарте, описывающем набор примитивов и операций разреженной линейной алгебры, предназначенных для построения алгоритмов анализа графов.
Он вводит абстракции моноидов и полуколец, объекты (скаляры, векторы, матрицы) и операции над ними (поэлементные, матричное умножение, маскирование, редукция).
Стандарт позволяет параметризовать операции пользовательскими полукольцами и тем самым единообразно выражать широкий класс графовых алгоритмов~--- от поиска кратчайших путей (тропическое полукольцо) до вычисления достижимости (булево полукольцо).

На практике этот стандарт реализуется в нескольких библиотеках: наиболее известная из них~--- SuiteSparse:GraphBLAS~\sidecite{Davis2018Algorithm9S}, широко используемая в высокопроизводительных графовых вычислениях.
Активно развиваются и GPGPU-реализации, такие как GraphBLAST.

Независимо от конкретной реализации, главное преимущество GraphBLAS-подхода для программиста~--- возможность выражать графовые алгоритмы через композицию матричных операций, делегируя заботу о разреженности, параллелизме и низкоуровневых оптимизациях библиотеке.
Операции $\texttt{map2}$ и $\texttt{mxm}$, параметризованные пользовательскими функциями, являются естественным обобщением этого подхода: они позволяют отказаться от жёсткого требования задавать полукольцо и вместо этого указывать ровно те операции, которые нужны для конкретной задачи, с явными гарантиями корректности на уровне типов~\sidecite{7761646}.
